%%%%%%%%%%%%%%%%%%%%%%%%%%%%%%%%%%%%%%%%%
% Medium Length Professional CV
% LaTeX Template
% Version 2.0 (8/5/13)
%
% This template has been downloaded from:
% http://www.LaTeXTemplates.com
%
% Original author:
% Trey Hunner (http://www.treyhunner.com/)
%
% Important note:
% This template requires the resume.cls file to be in the same directory as the
% .tex file. The resume.cls file provides the resume style used for structuring the
% document.
%
%%%%%%%%%%%%%%%%%%%%%%%%%%%%%%%%%%%%%%%%%

%----------------------------------------------------------------------------------------
%   PACKAGES AND OTHER DOCUMENT CONFIGURATIONS
%----------------------------------------------------------------------------------------

\documentclass{resume} % Use the custom resume.cls style
\usepackage[left=0.6in,top=0.4in,right=0.7in,bottom=0.5in]{geometry} % Document margins
\linespread{0.92}
\usepackage{hyperref}
\hypersetup{
    colorlinks=true,
    urlcolor=blue,
}

%----------------------------------------------------------------------------------------
%   DOCUMENT START
%----------------------------------------------------------------------------------------
\name{Chenming Ge} % Your name
\address{gecm@umich.edu \enspace\textbullet\enspace (734) 358-2718
\enspace\textbullet\enspace oscar-ge.github.io}

\begin{document}

%----------------------------------------------------------------------------------------
%   EDUCATION SECTION
%----------------------------------------------------------------------------------------


\begin{rSection}{Education}

\begin{rSubsection}{University of Michigan, Ann Arbor}{Ann Arbor, USA}{B.S.E. in Computer Science}{Aug 2025 - Present} 
\item \textbf{GPA:} 3.88/4.00
\item \textbf{Relevant Coursework:} \textbf{Machine Learning}, Intro to Algorithmic Robotics (A),  Operating Systems, Applied Parallel Programming with GPUs
\end{rSubsection}

\begin{rSubsection}{Shanghai Jiao Tong University}{Shanghai, China}{B.Eng. in Electrical and Computer Engineering}{Aug 2023 - Present}
\end{rSubsection}
\end{rSection}

\begin{rSection}{Research Interests}
AI Agents (with a focus on Web Agents and their memory), Robot Learning (VLA, VAM), and Efficient Machine Learning.
\end{rSection}

%----------------------------------------------------------------------------------------
%   RESEARCH EXPERIENCE 
%----------------------------------------------------------------------------------------

\begin{rSection}{Research Experience}

% Highlight your current independent research first
\begin{rSubsection}{Efficient Inference for Embodied Foundation Models}{Ann Arbor, USA}{Research Assistant, Advisor: Dr. Jiachen Liu}{Nov 2025 -- Present}
\item Profiled \textbf{COSMOS-Policy} across denoising steps; identified that block-level residual skipping is infeasible while cross-attention outputs remain highly stable (cosine $>$ 0.999).
\item Validated \textbf{cross-attention KV caching} on \textbf{24 RoboCasa tasks}: achieved  identical task success rate of 68.06\% to the baseline across 3,600 rollout trials with consistent denoising speedup across \textbf{all 24 tasks}.
\item Designing \textbf{task-aware token compression} that exploits the model's latent-frame slot structure to selectively prune image tokens via self-attention reduction, targeting further denoising acceleration on top of the validated caching strategy.
\end{rSubsection}

\begin{rSubsection}{Gravitational Effect on Swarming Behavior of Microorganisms}{Shanghai, China}{Research Assistant, Advisor: Prof. Zijie Qu}{Sep 2024 - Aug 2025}
\item Trained a \textbf{U-Net} model for colony boundary segmentation; diagnosed systematic failure modes, then adopted a \textbf{SAM}-based pipeline that achieved reliable automated detection, replacing manual annotation.
\end{rSubsection}


\end{rSection}

%----------------------------------------------------------------------------------------
%   PROJECTS 
%----------------------------------------------------------------------------------------

\begin{rSection}{Projects}
\begin{rSubsection}{Agent Memory Decay Under Controlled Web Drift -- \href{https://oscar-ge.github.io/files/webEvolve_progress_report.pdf}{[Progress Report]}}{Jan 2026 - Present}{EECS 545 (Machine Learning) Course Project, Advisor: Prof. Honglak Lee}{}
\item Developed an evaluation framework measuring Experience Transfer Gap (ETG) and Decay Rate (EDR) to study how test-time agent memory degrades under website interface evolution.
\item Created a reproducible testbed injecting 6 controlled GUI drift variants (e.g., structural, surface, content) to systematically benchmark failure modes.
\item Extracted and compared varying abstractions of test-time memory (e.g., semantic insights vs. procedural workflows) to quantify robustness against environmental coupling.
\end{rSubsection}

\begin{rSubsection}{Probabilistic Motion Planning for Redundant Robots}{Nov 2025 - Dec 2025}{EECS 465 (Intro to Algorithmic Robotics) Course Project, Advisor: Prof. Dmitry Berenson}{}
\item Reproduced and benchmarked \textbf{sampling-based motion planning algorithms} (RRT-Connect, PRM) for a \textbf{7-DOF Franka Panda} robot in simulation; implemented a hybrid sampling strategy that reduced trajectory generation latency by $\sim$75\%.
\end{rSubsection}

\end{rSection}

%----------------------------------------------------------------------------------------
%   TECHNICAL SKILLS SECTION
%----------------------------------------------------------------------------------------
\begin{rSection}{Skills}

\begin{tabular}{ @{} >{\bfseries}l @{\hspace{6ex}} l }
Languages & Mandarin (native), English (fluent, TOEFL 100) \\
Programming Languages & Python, C/C++, MATLAB, Java, LaTeX, Typst \\ 
Technical Skills & Hugging Face, PyTorch, Git, Claude Code. Linux, Docker, ssh, SolidWorks \\
\end{tabular}

\end{rSection}

\end{document}
